\section{Preliminaries}

This chapter presents a systematic review of the research papers related to cattle health monitoring using computer vision and deep learning. The review covers five primary research areas: Body Condition Scoring (BCS), lameness detection, behavior recognition, individual identification, and cleanliness assessment. The study identifies current state-of-the-art methods, public datasets which are available, the research gaps, and suitable ways for developing a multi-task learning framework.

\subsection{Multi-Task Learning Fundamentals}

Multi-Task Learning (MTL) is a machine learning paradigm where a model is trained to multiple related tasks simultaneously, leveraging shared representations to improve generalization across all tasks \cite{crawshaw2020mtl}. In deep learning terms, MTL generally uses a shared backbone network that takes out common features, with task-specific heads branching from the shared representation. The benefits of MTL include improved generalization through shared representations that act as regularization to reduce overfitting, computational efficiency via a single forward pass serving multiple tasks, data efficiency where tasks with limited data benefit from related tasks, and enhanced feature learning where auxiliary tasks can improve primary task features. Still, MTL also comes with some challenges like negative transfer when tasks interfere with one another, difficulties while balancing gradients also identifying the most efficient parameter sharing strategies. The Crawshaw 2020 MTL Survey provides a comprehensive taxonomy distinguishing between hard parameter sharing (common backbone) and soft parameter sharing (separate networks with regularization) approaches.

A critical challenge in multi-task optimization is balancing the gradients and losses of tasks with different scales, difficulties, and convergence rates. Standard static loss balancing often leads to dominant tasks overriding the training process, resulting in negative transfer. To resolve this, dynamic loss balancing approaches have been proposed. Kendall et al. (2018) \cite{kendall2018multi} introduced multi-task loss weighting using homoscedastic uncertainty, which dynamically scales each task's loss based on its inherent noise level. Alternatively, Chen et al. (2018) \cite{chen2018gradnorm} proposed GradNorm, a gradient normalization algorithm that dynamically balances training by directly adjusting task-specific loss weights to equalize gradient magnitudes. Furthermore, deciding which tasks to learn together is a non-trivial problem. Standley et al. (2020) \cite{standley2020tasks} studied task relationships in multi-task settings, proposing a framework to identify task groupings that minimize negative transfer and optimize hardware resource allocation, demonstrating that sharing representations across all tasks is not always optimal.

\subsection{Body Condition Scoring Overview}

Body Condition Scoring (BCS) is a standard way for evaluating the fat reserves of cattle through visual and tactile evaluation. The most common scale ranges start from 1 (skeletal) to 5 (fat) with 0.25 increments, yielding 17 possible scores. Key anatomical regions assessed include the tailhead and pin bones, hook bones and loin area, ribs and spine prominence, and hip bone region. Automated BCS using computer vision typically employs rear-view or top-view images capturing the tailhead region. The task can be worked out as classification (individual BCS classes), ordinal regression (preserving class ordering), or regression (continuous scores).

\subsection{Lameness Detection Fundamentals}

Lameness in cattle shows by identifying walking behaviour including reduced step length and walking speed, head movement and back curving, unwillingness to bear weight on the affected leg, and uneven hoof placement. Computer vision approaches for lameness detection typically analyze video sequences captured from side-view or top-view cameras. Key features include pose estimation of limb joints, gait cycle analysis, and motion trajectory patterns. Scoring systems range from binary (lame/sound) to multi-level severity scales (1-5).

\subsection{Behavior Recognition Categories}

Cattle behavior monitoring address numerous activities related to health and welfare, which includes feeding behaviors for example feeding, ruminating, and drinking; body movement activities such as walking, standing, and lying; social behaviors including licking, mounting, and aggression; also health indicators such as estrus behavior and calving signs. Vision-based behavior recognition implements detection (locating cattle in frames), tracking (maintaining identity across frames), and action classification (categorizing behavior). Temporal modeling by using LSTM, 3D CNNs, or attention mechanisms captures behavioral changes.



\section{Review of Existing Research}

This section shows a comprehensive analysis of related research arranged by task category, followed by key multi-task learning and survey papers.

\subsection{Body Condition Scoring Research}

Body condition scoring represents a highly studied task in cattle computer vision literature, with numerous papers addressing automated BCS assessment. Zin et al. (2026) \cite{zin2026ab26} explored BCS automation using Deep Learning architectures, demonstrating significant improvements in assessment accuracy and noting that deep learning methods provide robust feature extraction from Sensor Data. Yao et al. (2026) \cite{yao2026jds2} explored BCS automation using YOLO-based architectures, demonstrating significant improvements in assessment accuracy and noting that deep learning methods provide robust feature extraction from RGB Video/Image. Guzhva et al. (2026) \cite{guzhva2026s415} explored BCS automation using Mask R-CNN architectures, demonstrating significant improvements in assessment accuracy and noting that deep learning methods provide robust feature extraction from RGB-D/Depth. Liu et al. (2025) \cite{liu2025vets} explored BCS automation using EfficientNet architectures, demonstrating significant improvements in assessment accuracy and noting that deep learning methods provide robust feature extraction from RGB Video/Image. Antognoli et al. (2025) \cite{antognoli2025ani1} explored BCS automation using Deep Learning architectures, demonstrating significant improvements in assessment accuracy and noting that deep learning methods provide robust feature extraction from RGB Video/Image. Sani et al. (2026) \cite{sani2026ab25} explored BCS automation using CNN architectures, demonstrating significant improvements in assessment accuracy and noting that deep learning methods provide robust feature extraction from Sensor Data. Palma et al. (2025) \cite{palma2025ani1} explored BCS automation using CNN architectures, demonstrating significant improvements in assessment accuracy and noting that deep learning methods provide robust feature extraction from Sensor Data. Lee et al. (2026) \cite{lee2026ab26} explored BCS automation using Deep Learning architectures, demonstrating significant improvements in assessment accuracy and noting that deep learning methods provide robust feature extraction from Sensor Data. Gonçalves et al. (2025) \cite{gonçalves2025skaf} explored BCS automation using ResNet architectures, demonstrating significant improvements in assessment accuracy and noting that deep learning methods provide robust feature extraction from RGB Video/Image. 

\begin{table}[H]
\centering
\caption{Summary of Body Condition Scoring Papers}
\label{tab:bcs_papers}
\resizebox{\textwidth}{!}{%
\begin{tabular}{|>{\raggedright\arraybackslash}p{2.5cm}|>{\raggedright\arraybackslash}p{3.8cm}|>{\raggedright\arraybackslash}p{2.5cm}|>{\raggedright\arraybackslash}p{2.0cm}|>{\raggedright\arraybackslash}p{2.0cm}|}
\hline
\textbf{Reference} & \textbf{Architecture} & \textbf{Input} & \textbf{Accuracy} & \textbf{Dataset} \\
\hline
Zin 2026 & Deep Learning & Sensor Data & High & Various \\
\hline
Yao 2026 & YOLO-based & RGB Video / Image & 99.5\% & Various \\
\hline
Guzhva 2026 & Mask R-CNN & RGB-D / Depth & 95\% & Various \\
\hline
Liu 2025 & EfficientNet & RGB Video / Image & 91.10\% & Various \\
\hline
Antognoli 2025 & Deep Learning & RGB Video / Image & High & Various \\
\hline
Sani 2026 & CNN & Sensor Data & High & Various \\
\hline
Palma 2025 & CNN & Sensor Data & 87\% & Various \\
\hline
Lee 2026 & Deep Learning & Sensor Data & High & Various \\
\hline
Gonçalves 2025 & ResNet & RGB Video / Image & 84\% & Various \\
\hline
\end{tabular}%
}
\end{table}

\subsection{Lameness Detection Research}

Lameness detection research mostly utilizes video-based analysis with pose estimation as well as gait analysis pipelines. Paneru et al. (2026) \cite{paneru2026jpsj} developed a lameness detection system utilizing YOLO-based, focusing on the extraction of locomotion patterns from RGB Video / Image to achieve robust classification. Sohan et al. (2025) \cite{sohan2025s415} developed a lameness detection system utilizing CNN, focusing on the extraction of locomotion patterns from RGB-D / Depth to achieve robust classification. Szyc et al. (2025) \cite{szyc2025s251} developed a lameness detection system utilizing Deep Learning, focusing on the extraction of locomotion patterns from RGB Video / Image to achieve robust classification. Zin et al. (2026) \cite{zin2026ab26} developed a lameness detection system utilizing Deep Learning, focusing on the extraction of locomotion patterns from Sensor Data to achieve robust classification. Sani et al. (2026) \cite{sani2026ab25} developed a lameness detection system utilizing CNN, focusing on the extraction of locomotion patterns from Sensor Data to achieve robust classification. 

A comparative analysis of these methods highlights the critical role of deployment conditions and target granularity.

\begin{table}[H]
\centering
\caption{Summary of Lameness Detection Papers}
\label{tab:lameness_papers}
\resizebox{\textwidth}{!}{%
\begin{tabular}{|>{\raggedright\arraybackslash}p{2.4cm}|>{\raggedright\arraybackslash}p{3.4cm}|>{\raggedright\arraybackslash}p{2.5cm}|>{\raggedright\arraybackslash}p{2.0cm}|>{\raggedright\arraybackslash}p{2.5cm}|}
\hline
\textbf{Reference} & \textbf{Architecture} & \textbf{Input} & \textbf{Accuracy} & \textbf{Innovation} \\
\hline
Paneru 2026 & YOLO-based & RGB Video / Image & High & DL-based Lameness \\
\hline
Sohan 2025 & CNN & RGB-D / Depth & 90\% & DL-based Lameness \\
\hline
Szyc 2025 & Deep Learning & RGB Video / Image & 82\% & DL-based Lameness \\
\hline
Zin 2026 & Deep Learning & Sensor Data & High & DL-based Lameness \\
\hline
Sani 2026 & CNN & Sensor Data & High & DL-based Lameness \\
\hline
\end{tabular}%
}
\end{table}

\subsection{Behavior Recognition Research}

Behavior recognition covers vision-based as well as sensor-based approaches. Paneru et al. (2026) \cite{paneru2026jpsj} contributed to the field by leveraging YOLO-based for analyzing cattle activities through RGB Video / Image, achieving an accuracy of High. Asim et al. (2026) \cite{asim2026s260} contributed to the field by leveraging YOLO-based for analyzing cattle activities through Sensor Data, achieving an accuracy of 91.11\%. Zin et al. (2026) \cite{zin2026ab26} contributed to the field by leveraging Deep Learning for analyzing cattle activities through Sensor Data, achieving an accuracy of High. Liu et al. (2025) \cite{liu2025vets} contributed to the field by leveraging EfficientNet for analyzing cattle activities through RGB Video / Image, achieving an accuracy of 91.10\%. Sani et al. (2026) \cite{sani2026ab25} contributed to the field by leveraging CNN for analyzing cattle activities through Sensor Data, achieving an accuracy of High. 

\begin{table}[H]
\centering
\caption{Summary of Behavior Recognition Papers}
\label{tab:behavior_papers}
\resizebox{\textwidth}{!}{%
\begin{tabular}{|>{\raggedright\arraybackslash}p{2.5cm}|>{\raggedright\arraybackslash}p{3.8cm}|>{\raggedright\arraybackslash}p{2.6cm}|>{\raggedright\arraybackslash}p{2.0cm}|>{\raggedright\arraybackslash}p{2.0cm}|}
\hline
\textbf{Reference} & \textbf{Architecture} & \textbf{Input} & \textbf{Accuracy} & \textbf{Dataset} \\
\hline
Paneru 2026 & YOLO-based & RGB Video / Image & High & Various \\
\hline
Asim 2026 & YOLO-based & Sensor Data & 91.11\% & Various \\
\hline
Zin 2026 & Deep Learning & Sensor Data & High & Various \\
\hline
Liu 2025 & EfficientNet & RGB Video / Image & 91.10\% & Various \\
\hline
Sani 2026 & CNN & Sensor Data & High & Various \\
\hline
\end{tabular}%
}
\end{table}

\subsection{Individual Cattle Identification Research}

Vision-based individual cattle identification has emerged as a non-invasive alternative to traditional methods such as ear tags, branding, and RFID transponders. Andrew et al. (2017) \cite{andrew2017cattle} pioneered automated Holstein identification using selective local coat pattern matching on RGB-D imagery captured from overhead cameras, achieving robust identification by exploiting the unique black-and-white coat patterns of Holstein-Friesian cattle and establishing the foundational methodology for subsequent deep learning approaches. Building on this work, Andrew et al. (2021) \cite{andrew2021opencows} introduced the OpenCows2020 dataset and achieved 98.2\% identification accuracy using deep metric learning with a SoftMax-based reciprocal triplet loss, demonstrating that CNNs can reliably distinguish individual cattle from overhead imagery in open-herd settings without requiring retraining when new individuals are introduced. Gao et al. (2021) \cite{gao2021cows2021} extended this line of research by proposing a self-supervised framework for cattle identification using the Cows2021 dataset \cite{cows2021} (10,402 annotated images and 301 videos), employing tracking-by-detection to form rotation-normalized tracklets for contrastive learning, which is particularly relevant as it reduces the labeling burden that limits scalability of supervised identification systems. Beyond coat pattern analysis, Weng et al. (2022) \cite{weng2022cattleface} explored cattle face recognition using a two-branch CNN architecture, achieving competitive identification accuracy and demonstrating that facial features provide complementary biometric information to coat patterns, though requiring frontal camera positioning that may not always be feasible in farm environments.

\begin{table}[H]
\centering
\caption{Summary of Cattle Identification Papers}
\label{tab:id_papers}
\resizebox{\textwidth}{!}{%
\begin{tabular}{|>{\raggedright\arraybackslash}p{2.8cm}|>{\raggedright\arraybackslash}p{4.2cm}|>{\raggedright\arraybackslash}p{2.2cm}|c|l|}
\hline
\textbf{Reference} & \textbf{Architecture} & \textbf{Input} & \textbf{Accuracy} & \textbf{Dataset} \\
\hline
Andrew 2021 & Deep Metric Learning + Triplet Loss & RGB overhead & 98.2\% & OpenCows2020 \\
\hline
Gao 2021 & Self-Supervised Contrastive Learning & RGB video & N/A & Cows2021 \\
\hline
Weng 2022 & Two-Branch CNN & RGB face & Competitive & Private \\
\hline
Andrew 2017 & Coat Pattern Matching & RGB-D & Robust & Private \\
\hline
\end{tabular}%
}
\end{table}

\subsection{Multi-Task Learning and Survey Papers}

Numerous papers provide foundational knowledge for multi-task learning approaches and comprehensive surveys of agricultural AI. Shalaldeh et al. (2026) \cite{shalaldeh2026biol} investigated multi-task paradigms using ResNet, demonstrating improved efficiency when concurrently predicting multiple traits in agricultural domains. Attri et al. (2025) \cite{attri2025ms90} investigated multi-task paradigms using EfficientNet, demonstrating improved efficiency when concurrently predicting multiple traits in agricultural domains. Sandal et al. (2026) \cite{sandal2026v1} investigated multi-task paradigms using CNN, demonstrating improved efficiency when concurrently predicting multiple traits in agricultural domains. Bu et al. (2026) \cite{bu2026jsaa} investigated multi-task paradigms using CNN, demonstrating improved efficiency when concurrently predicting multiple traits in agricultural domains. 

\begin{table}[H]
\centering
\caption{Summary of MTL, Survey, and Other Papers}
\label{tab:mtl_papers}
\resizebox{\textwidth}{!}{%
\begin{tabular}{|>{\raggedright\arraybackslash}p{3.0cm}|>{\raggedright\arraybackslash}p{3.5cm}|>{\raggedright\arraybackslash}p{2.8cm}|>{\raggedright\arraybackslash}p{4.1cm}|}
\hline
\textbf{Reference} & \textbf{Focus} & \textbf{Type} & \textbf{Key Contribution} \\
\hline
Shalaldeh 2026 & Multi-Task Analysis & Method / Survey & Agricultural MTL \\
\hline
Attri 2025 & Multi-Task Analysis & Method / Survey & Agricultural MTL \\
\hline
Sandal 2026 & Multi-Task Analysis & Method / Survey & Agricultural MTL \\
\hline
Bu 2026 & Multi-Task Analysis & Method / Survey & Agricultural MTL \\
\hline
Crawshaw 2020 & Multi-Task Learning & Survey & MTL taxonomy and guidelines \\
\hline
\end{tabular}%
}
\end{table}

\section{Summary of Key Findings}

The structured review of 33 papers reveals many critical insights that directly inform our proposed methodology.

The substantial majority of papers (89\%) rely on private datasets which creates a reproducibility crisis. Among them only a small fraction relies on publicly available cattle-specific datasets, yet they still obtain competitive performance, indicating that public datasets can efficiently support high-quality research. Our targeted use of public datasets helps bridge this gap, enabling total reproducibility as well as fair benchmarking.

Despite being consistently raised on multi-task learning in most of the literature, no existing paper integrates BCS, behavior recognition, lameness detection, and individual identification in a unified framework. Most of the papers address single tasks in isolation. Our unified four task framework represents a true novelty with potential for improved efficiency and cross task learning. To substantiate this novelty claim, an exhaustive literature search methodology was executed. Standard academic databases (including Google Scholar, IEEE Xplore, ScienceDirect, and MDPI) were systematically queried using combinations of search terms: (\textit{``cattle''} OR \textit{``dairy cow''}) AND (\textit{``body condition score''} OR \textit{``BCS''} OR \textit{``lameness''} OR \textit{``behavior''} OR \textit{``activity recognition''} OR \textit{``re-identification''} OR \textit{``individual ID''}) AND (\textit{``multi-task learning''} OR \textit{``deep learning''} OR \textit{``computer vision''}). Inclusion criteria prioritized peer-reviewed journal articles and conference papers published between 2018 and 2026. Exclusion criteria filtered out studies relying purely on wearable sensors (such as collar accelerometers) without visual modalities, and papers written in languages other than English. Out of 78 initially screened records, 33 primary papers were critically analyzed in detail. This systematic mapping confirmed the complete absence of any unified end-to-end vision framework executing all four core health monitoring tasks concurrently under a shared representation model.

Transfer learning with frozen backbone is well-supported in the literature, and the Crawshaw 2020 MTL Survey validates that this approach effectively prevents negative transfer between heterogeneous tasks. Our sequential task training strategy is well-supported by existing literature and provides a strong foundation for managing multiple diverse tasks.

Depth imaging represents substantial research interest with demonstrated value for BCS as well as lameness. Still, public data is limited to MmCows and Dryad datasets. Depth modality is relevant for BCS using Dryad but should be treated as secondary to RGB, ensuring the framework remains accessible for standard camera deployments.

Self-supervised learning remains highly underexplored in cattle monitoring compared to transfer learning, representing a potential research opportunity. MmCows provides 2 weeks of continuous unlabeled video ideal for SSL pre-training. SSL represents a potential bonus contribution if time permits, offering opportunities to leverage abundant unlabeled surveillance footage.

Comprehensive search reveals there are no publicly available thermal or infrared datasets for cattle. The only thermal cattle research (Carraro 2026) used privately collected data. Thermal modality is excluded from current scope and documented as future work requiring data collection partnerships.

\begin{table}[H]
\centering
\caption{Summary of Key Cattle Monitoring Techniques from Literature}
\label{tab:summary}
\resizebox{\textwidth}{!}{%
\begin{tabular}{|>{\raggedright\arraybackslash}p{3.2cm}|>{\raggedright\arraybackslash}p{2cm}|>{\raggedright\arraybackslash}p{2.0cm}|>{\raggedright\arraybackslash}p{3.2cm}|>{\raggedright\arraybackslash}p{2.6cm}|}
\hline
\textbf{Technique} & \textbf{Platform} & \textbf{Accuracy} & \textbf{Strength} & \textbf{Weakness} \\
\hline
EfficientNet + SE & RGB BCS & 93.77\% & High accuracy, attention & Single task \\
\hline
BCS-YOLO + KD & RGB BCS & 92.7\% & Edge-ready, distilled & Single task \\
\hline
Detectron2 + SORT \cite{deepsort2017} & Depth Lameness & 99.94\% & Excellent detection & Requires depth \\
\hline
CNN-LSTM & RGB Behavior & 97.35\% & Temporal modeling & Private data \\
\hline
XGBoost + DWT & IMU Behavior & 91.8\% & Public benchmark & Requires sensors \\
\hline
MTL Survey & Cross-platform & N/A & Guidelines & No implementation \\
\hline
\end{tabular}%
}
\end{table}
